Booked-load Dijkstra is ahead of battery Dijkstra by 0.4 to 44.3 points; the 95\% interval is above zero in 4 of 5 conditions but includes zero for $n=50$. The two coordination mechanisms are not interchangeable once the operating point changes. Flow augmentation is ahead for $n=50$ ($-3.6$) and clustered layouts ($-4.0$); booked-load Dijkstra is ahead for $n=200$ ($+0.9$), the corner sink ($+6.3$), clustered layouts with a corner sink ($+25.5$). Booked load gains most where traffic funnels hardest towards the sink, while flow augmentation does better in small or clustered networks. The LP planner is best wherever it was run except for clustered layouts with a corner sink, where booked-load Dijkstra reaches 59.8\% against 49.8\%. A plausible cause is that, with very long links, rounding the fractional flow to one next hop per node per round loses more than the planner gains. Absolute scores are also low there for every router (the best reaches about 60\%), so the bound itself may be loose in that condition.
